% =============================================================================
%  §3  Research Methodology — Systematic Literature Review
%
%  Governing frameworks (read these before writing each paragraph):
%    - Xiao & Watson (2019) — 3-phase, 8-step SLR process
%        Phase 1 Planning:   Step 1 Formulate the problem
%                            Step 2 Develop and validate the protocol
%        Phase 2 Conducting: Step 3 Search the literature
%                            Step 4 Screen for inclusion
%                            Step 5 Assess quality
%                            Step 6 Extract data
%                            Step 7 Analyse and synthesise
%        Phase 3 Reporting:  Step 8 Report findings
%    - Wohlin (2014)        — Snowballing: start set → backward → forward
%                             → author contact → data extraction
%    - Page et al. (2021)   — PRISMA 2020 flow and reporting checklist
%    - Kitchenham & Charters (2007) — protocol structure and quality criteria
%
%  WRITING RULES:
%    - Every subsection maps to one or more of the 8 steps above.
%      % REQUIRED BY comments tell you which step and what it demands.
%    - Write your OWN prose for each paragraph.  The bullet hints under
%      each %% WRITE: comment tell you exactly what content must appear.
%    - \TODO{...} = a number or phrase to insert after the pipeline runs.
%    - Do NOT remove any subsection or table — they are all required by
%      the frameworks.  You may add subparagraphs but not delete structure.
% =============================================================================

% Helper command — remove before submission:
\newcommand{\TODO}[1]{\textbf{[TODO: #1]}}

\section{Research Methodology}
\label{sec:methodology}

% =============================================================================
\subsection{Review Type and Governing Framework}
\label{sec:methodology:framework}
% =============================================================================

% REQUIRED BY: Xiao & Watson (2019) Step 1 — Formulate the problem
%              Xiao & Watson (2019) Step 2 — Develop and validate the protocol
%              Kitchenham & Charters (2007) — SLR definition and pre-specification
%              Page et al. (2021) — PRISMA Methods: describe protocol

%% WRITE ¶1 — SLR definition and PRISMA framework (~3–4 sentences):
%   • Name the methodology: Systematic Literature Review.
%   • Give the standard definition (pre-planned, comprehensive, reproducible).
%   • State which reporting standard governs the review (PRISMA 2020).
%   • Name the four PRISMA stages: Identification, Screening, Eligibility,
%     Inclusion.
%   • Cite: \citep{XiaoWatson2019, KitchenhamCharters2007, Page2021}

\bigskip\noindent\textit{[¶1 — SLR definition and PRISMA framework]}
% ↑ Delete this line and write your paragraph here.

%% WRITE ¶2 — Research question using PICOC decomposition (~3–5 sentences):
%   REQUIRED BY: Xiao & Watson (2019) Step 1 — research question must be
%                stated explicitly; PICO/PICOC decomposition recommended.
%   • State the research question in full (blockquote format is conventional):
%       "What methods have been proposed for decentralised, adapter-based
%        multi-task inference over a shared frozen transformer backbone,
%        and what are the open research gaps?"
%   • Decompose it using PICOC (Population / Intervention / Comparison /
%     Outcome / Context):
%       P — large language models deployed in distributed / P2P environments
%       I — parameter-efficient fine-tuning (PEFT) and adapter modules
%       C — centralised fine-tuning and monolithic serving baselines
%       O — multi-task inference capability, serving efficiency, modularity
%       C — peer-to-peer networks; resource-constrained inference
%   • State the three intersecting fields the review synthesises:
%     (i) PEFT / adapter architectures; (ii) distributed / P2P ML systems;
%     (iii) multi-task LLM serving and inference.
%   • Cite: \citep{XiaoWatson2019, KitchenhamCharters2007}

\bigskip\noindent\textit{[¶2 — Research question and PICOC decomposition]}
% ↑ Delete this line and write your paragraph here.

\begin{quote}
  \textit{What methods have been proposed for decentralised, adapter-based
  multi-task inference over a shared frozen transformer backbone, and what
  are the open research gaps?}
\end{quote}

%% WRITE ¶3 — Protocol pre-specification and audit trail (~2–3 sentences):
%   REQUIRED BY: Xiao & Watson (2019) Step 2 — a protocol must be written
%                BEFORE searching begins; any amendment must be logged.
%                Kitchenham & Charters (2007) — criteria defined a priori.
%   • State that the review protocol was defined prior to execution.
%   • Name what the protocol documents: research question, I/E criteria,
%     search strategy, screening procedure, enrichment pipeline, and
%     data extraction schema.
%   • Note where it is archived (repository) and how changes are tracked
%     (SHA-256 checksums, version-controlled CSVs and JSON logs).
%   • Cite: \citep{XiaoWatson2019, KitchenhamCharters2007}

\bigskip\noindent\textit{[¶3 — Protocol pre-specification and audit trail]}
% ↑ Delete this line and write your paragraph here.
%
% Footnote suggestion:
%   \footnote{Protocol files: \texttt{SLR\_methodology\_outline.md},
%             \texttt{app/screen.py}, \texttt{app/enrich\_and\_filter.py};
%             SHA-256 checksums retained per intermediate output file.}


% =============================================================================
\subsection{Seed Paper Selection and Start Set Validation}
\label{sec:methodology:seeds}
% =============================================================================

% REQUIRED BY: Wohlin (2014) §3.1 — Start Set construction and validation
%              Wohlin (2014) Lesson: start set must cover multiple communities
%              and be validated via full-paper reading before snowballing.
%              Xiao & Watson (2019) Step 3 — justify why these sources.

%% WRITE ¶4 — Start set selection criteria and characteristics (~4–6 sentences):
%   • Introduce the start set (G0, 9 papers) as the snowballing anchor.
%   • State the three selection criteria (Wohlin 2014 requires justifying
%     each characteristic):
%       (i)  Relevance to ≥ 2 of the 3 core themes
%       (ii) Citation count ≥ 300 at time of selection
%       (iii) Publication in a top-tier venue
%   • Address Wohlin's 5 required characteristics of a good start set:
%       – Community coverage: seeds span adapters, P2P systems, LLM serving
%         (avoids isolated citation clusters)
%       – Size: 9 papers appropriate to the breadth of a 3-theme review
%       – Diversity: multiple publishers, years, and author groups
%       – Synonym coverage: seeds use varied terminology (LoRA, adapters,
%         PEFT, collaborative inference, multi-tenant serving)
%       – Drawn from key terms in the research question
%   • List the 9 seeds briefly (cite each):
%     \citep{Houlsby2019, Hu2022, Pfeiffer2020hub, Pfeiffer2021fusion,
%            Sajina2024, Borzunov2023, Han2024, Sheng2024}
%   • Cite: \citep{Wohlin2014}

\bigskip\noindent\textit{[¶4 — Start set selection criteria and Wohlin's 5 characteristics]}
% ↑ Delete this line and write your paragraph here.

%% WRITE ¶5 — Start set validation (~2–3 sentences):
%   REQUIRED BY: Wohlin (2014) §3.1 — "the full paper should be examined
%                before deciding to use it as a paper in the snowballing."
%                Papers not validated cannot enter snowballing without risk
%                of rollback.
%   • State that each G0 seed was confirmed against the full inclusion criteria
%     on the basis of full-text reading before being used as a snowball anchor.
%   • Note that seeds lacking a resolvable API identifier (DOI or arXiv ID)
%     were excluded from programmatic snowballing.
%   • Cite: \citep{Wohlin2014}

\bigskip\noindent\textit{[¶5 — Start set validation (full-text confirmation)]}
% ↑ Delete this line and write your paragraph here.


% =============================================================================
\subsection{Pre-Validated Corpora (Groups G1--G6)}
\label{sec:methodology:prevalidated}
% =============================================================================

% REQUIRED BY: Xiao & Watson (2019) Step 3 — document all sources searched,
%              including venue-specific and hand-search sources.
%              Wohlin (2014) — complementary manual searches of key journals
%              and conferences are recommended.

%% WRITE ¶6 — Purpose and construction of G1–G6 (~4–5 sentences):
%   • Introduce the six pre-validated corpora as a complementary hand-search
%     of leading publication venues (Wohlin 2014 recommends searching specific
%     journals and conferences identified as relevant).
%   • State the construction method: targeted manual curation via
%     Undermind AI-assisted search against the same leading venues.
%   • Explain why they bypass the automated year and domain filters:
%     their relevance was established during curation (human-validated,
%     topic-specific queries); re-assessing programmatically would be
%     redundant.  The manual curation IS the quality gate.
%   • Reference Table~\ref{tab:prevalidated}.
%   • Cite: \citep{Wohlin2014, XiaoWatson2019}

\bigskip\noindent\textit{[¶6 — Pre-validated corpora: purpose and construction]}
% ↑ Delete this line and write your paragraph here.

\begin{table}[htbp]
  \centering
  \caption{Pre-validated corpora (Groups G1--G6) assembled by targeted
           manual curation of leading publication venues.}
  \label{tab:prevalidated}
  \begin{tabular}{llr}
    \toprule
    \textbf{Group} & \textbf{Theme} & \textbf{\textit{n}} \\
    \midrule
    G1 & PEFT methods beyond adapters and LoRA                   &  62 \\
    G2 & Adapter composition for multi-task NLP                  &  46 \\
    G3 & Decentralised and P2P machine learning systems          & 112 \\
    G4 & Adapter multiplexing for efficient LLM inference        &  64 \\
    G5 & Routing and mixture-of-experts for modular PEFT         &  69 \\
    G6 & Federated PEFT for transformer NLP                      &  91 \\
    \midrule
    \textbf{Total} & & \textbf{444} \\
    \bottomrule
  \end{tabular}
\end{table}


% =============================================================================
\subsection{Citation Snowballing Procedure}
\label{sec:methodology:snowballing}
% =============================================================================

% REQUIRED BY: Wohlin (2014) §3 — full backward + forward procedure
%              Wohlin (2014) §3.2.1 — backward: 4-layer sequential screening
%              Wohlin (2014) §3.2.2 — forward: same layered logic
%              Wohlin (2014) §3.3  — author contact after saturation
%              Wohlin (2014) §3.4  — data extraction after inclusion
%              Xiao & Watson (2019) Step 3 — report search engines and date

% -----------------------------------------------------------------------
\subsubsection{Backward Snowballing}
\label{sec:methodology:snowballing:backward}
% -----------------------------------------------------------------------

%% WRITE ¶7 — Backward snowballing procedure (~4–5 sentences):
%   REQUIRED BY: Wohlin (2014) §3.2.1 — sequential 4-layer screening.
%   • Describe the purpose: reference lists of G0 seeds examined to surface
%     earlier work that influenced the seeds.
%   • Describe the 4-layer sequential screening applied to each reference:
%       Layer 1 — Basic exclusions directly from the reference list
%                 (no retrieval): language, year, document type.
%                 Exclude already-examined papers.
%       Layer 2 — Title, publication venue, and in-text citation context
%                 (information available within the citing paper, no
%                 retrieval needed).  Assess where and how the candidate
%                 is cited; a citation in a key related-work sentence
%                 is a strong positive signal.
%       Layer 3 — Abstract (retrieved if Layers 1–2 are inconclusive).
%       Layer 4 — Full text (retrieved only if abstract is inconclusive;
%                 required for a definitive inclusion decision).
%   • State that only papers passing all applicable layers were flagged
%     for inclusion.
%   • Cite: \citep{Wohlin2014}

\bigskip\noindent\textit{[¶7 — Backward snowballing: 4-layer sequential screening]}
% ↑ Delete this line and write your paragraph here.

% -----------------------------------------------------------------------
\subsubsection{Forward Snowballing}
\label{sec:methodology:snowballing:forward}
% -----------------------------------------------------------------------

%% WRITE ¶8 — Forward snowballing procedure (~3–4 sentences):
%   REQUIRED BY: Wohlin (2014) §3.2.2 — forward via citation index.
%   • Describe the purpose: identifying subsequent works that cite each G0
%     seed, capturing more recent developments.
%   • State the tool used: Semantic Scholar Academic Graph API
%     \citep{SemanticScholar} (primary source for forward citations).
%   • Describe the same 4-layer sequential screening (title/snippet →
%     abstract → context of citation → full text) applied to each citing
%     work, per Wohlin (2014).
%   • Cite: \citep{Wohlin2014, SemanticScholar}

\bigskip\noindent\textit{[¶8 — Forward snowballing: citation index and layered screening]}
% ↑ Delete this line and write your paragraph here.

% -----------------------------------------------------------------------
\subsubsection{Search Engines, Supplementary Sources, and Search Date}
\label{sec:methodology:snowballing:engines}
% -----------------------------------------------------------------------

%% WRITE ¶9 — Sources used and WoS role (~4–5 sentences):
%   REQUIRED BY: Xiao & Watson (2019) Step 3 — list every database/source;
%                report exact search date; no single database is sufficient.
%                Wohlin (2014) — Semantic Scholar recommended for forward;
%                complementary manual search of key venues recommended.
%   • Name the primary snowball engine: Semantic Scholar Academic Graph API.
%   • Name the supplementary engines: Scopus and ACL Anthology
%     (capture venue-specific corpora not fully indexed by SS).
%   • Clarify the WoS role: post-hoc enrichment only — appending WoS
%     accession numbers and citation counts to already-retrieved records;
%     NOT a source of new snowball candidates.
%   • State the search execution date (from pipeline log):
%     \TODO{insert date from \texttt{log\_retrieval\_*.json} field
%     \texttt{timestamp\_utc}}.
%   • Cite: \citep{SemanticScholar, XiaoWatson2019}

\bigskip\noindent\textit{[¶9 — Search engines, WoS role, and search date]}
% ↑ Delete this line and write your paragraph here.

% -----------------------------------------------------------------------
\subsubsection{Deduplication, Saturation Assessment, and Iteration Scope}
\label{sec:methodology:snowballing:saturation}
% -----------------------------------------------------------------------

%% WRITE ¶10 — Deduplication and raw counts (~2–3 sentences):
%   REQUIRED BY: Wohlin (2014) — duplicates removed before screening;
%                report per-iteration counts.
%                Page et al. (2021) PRISMA — records retrieved and
%                duplicates removed must be reported.
%   • State total raw records retrieved: 1,150.
%   • State deduplication method: by DOI and Semantic Scholar paper ID.
%   • State result: 41 unique new records (1,109 duplicates removed).
%   • Cite: \citep{Wohlin2014, Page2021}

\bigskip\noindent\textit{[¶10 — Deduplication: 1,150 raw → 41 unique new]}
% ↑ Delete this line and write your paragraph here.

%% WRITE ¶11 — Iteration scope and saturation (~3–4 sentences):
%   REQUIRED BY: Wohlin (2014) — stopping criterion is "no new papers
%                found in an iteration of both backward and forward";
%                Lesson 7: inclusion frequency should decrease each
%                iteration; Lesson 8: if not decreasing, search synonyms.
%   HONEST LIMITATION TO STATE:
%   • This review executed ONE wave of snowballing (not iterative
%     to zero).
%   • Justify: saturation was assessed MANUALLY rather than by iteration
%     exhaustion — the high overlap between newly retrieved snowball
%     candidates and the pre-validated G1–G6 corpora indicated that
%     additional waves would yield diminishing returns (consistent with
%     Wohlin's Lesson 7).
%   • Acknowledge this as a limitation and that the G1–G6 corpora
%     serve as a complementary safeguard against missed clusters
%     (Wohlin 2014: start set should cover multiple communities).
%   • Cite: \citep{Wohlin2014}
%
%   Per-iteration tracking table (fill in after pipeline run):

\bigskip\noindent\textit{[¶11 — Iteration scope, saturation justification, and limitation]}
% ↑ Delete this line and write your paragraph here.

% Iteration tracking table — REQUIRED BY Wohlin (2014) for traceability
\begin{table}[htbp]
  \centering
  \caption{Snowballing iteration summary (one wave executed).
           ``Candidates examined'' excludes immediately duplicated records.}
  \label{tab:iterations}
  \begin{tabular}{lrrr}
    \toprule
    \textbf{Direction} & \textbf{Candidates examined} &
    \textbf{Included} & \textbf{Cumulative included} \\
    \midrule
    Backward (Wave 1) & \TODO{$n$} & \TODO{$n$} & \TODO{$n$} \\
    Forward  (Wave 1) & \TODO{$n$} & \TODO{$n$} & \TODO{$n$} \\
    \midrule
    \textbf{Total (after deduplication)} & \TODO{1{,}150} &
    \textbf{41} & \textbf{41} \\
    \bottomrule
  \end{tabular}
\end{table}

%% NOTE on author contact (Wohlin 2014 §3.3):
%   Wohlin recommends contacting the authors of included papers after the
%   snowballing wave to surface grey literature, preprints, and papers in
%   obscure venues.  If you conducted this step, add a paragraph here.
%   If not, acknowledge it as a limitation in the limitations section.


% =============================================================================
\subsection{Inclusion and Exclusion Criteria}
\label{sec:methodology:criteria}
% =============================================================================

% REQUIRED BY: Xiao & Watson (2019) Step 4 — I/E criteria must be
%              pre-specified in the protocol; rationale for each stated.
%              Kitchenham & Charters (2007) — criteria defined a priori.
%              Page et al. (2021) — PRISMA §Methods: report eligibility criteria.

%% WRITE ¶12 — Introduction to criteria (~2 sentences):
%   • State that criteria were pre-specified in the protocol before
%     searching began (Xiao & Watson Step 2 requirement).
%   • State they were applied consistently across all screening stages.
%   • Reference Table~\ref{tab:ie-criteria}.
%   • Cite: \citep{XiaoWatson2019, KitchenhamCharters2007}

\bigskip\noindent\textit{[¶12 — Introduction to pre-specified criteria]}
% ↑ Delete this line and write your paragraph here.

\begin{table}[htbp]
  \centering
  \caption{Pre-specified inclusion and exclusion criteria (applied
           consistently across all screening stages).}
  \label{tab:ie-criteria}
  \begin{tabular}{p{3.0cm}p{5.4cm}p{5.4cm}}
    \toprule
    \textbf{Criterion} & \textbf{Inclusion} & \textbf{Exclusion} \\
    \midrule
    Publication year
      & $\geq 2021$ (snowballed records; pre-2021 foundational works
        enter via G0--G6 manual curation)
      & $< 2021$ \\[4pt]
    Language
      & English (implicit; source databases index predominantly
        English-language works; no mechanical filter applied)
      & Non-English works not surfaced by the APIs \\[4pt]
    Document type
      & Peer-reviewed conference paper, journal article, or arXiv
        preprint with $\geq 1$ citation
      & Blog posts, grey literature, software documentation, workshop
        papers without formal proceedings \\[4pt]
    Minimum topical relevance
      & Addresses $\geq 1$ of: (a)~PEFT / adapters; (b)~distributed /
        P2P / federated ML; (c)~multi-task LLM serving or inference
      & Exclusively covers unrelated domains (CV, audio, RL) with no
        NLP or systems angle \\[4pt]
    Tier~1 priority
      & Addresses all three themes jointly: PEFT + distributed/P2P +
        LLM serving
      & --- \\[4pt]
    Abstract availability
      & Abstract accessible via API or pre-validated corpus
      & No accessible metadata and no resolvable DOI \\
    \bottomrule
  \end{tabular}
\end{table}

%% WRITE ¶13 — Year cutoff rationale (~2–3 sentences):
%   REQUIRED BY: Xiao & Watson (2019) — rationale for every criterion.
%   • Explain why 2021 was chosen as the lower bound for snowballed records.
%   • Key argument: LoRA \citep{Hu2022} and widespread large-scale
%     transformer deployment mark the practical onset of the PEFT-at-scale
%     era this thesis addresses.
%   • Note that pre-2021 foundational works (e.g., \citealt{Houlsby2019})
%     are included via the manually curated G0--G6 corpora which bypass the
%     year filter.

\bigskip\noindent\textit{[¶13 — Year cutoff rationale]}
% ↑ Delete this line and write your paragraph here.


% =============================================================================
\subsection{Screening Procedure}
\label{sec:methodology:screening}
% =============================================================================

% REQUIRED BY: Xiao & Watson (2019) Step 4 — three-stage sequential
%              screening: title → abstract → full text; inter-rater
%              reliability required.
%              Wohlin (2014) §3.2 — same sequential logic within snowballing.
%              Page et al. (2021) — PRISMA Screening stage flow.

%% WRITE ¶14 — Overview of title screening (~2–3 sentences):
%   • State total records entering the pipeline: 1,008
%     (41 snowballed + 444 pre-validated + \TODO{any others}).
%   • Describe the three-layer approach: sequential decision hierarchy
%     following Wohlin (2014) and the coarse-to-fine sieve of
%     Xiao & Watson (2019).
%   • Reference Table~\ref{tab:screening}.
%   • Cite: \citep{Wohlin2014, XiaoWatson2019}

\bigskip\noindent\textit{[¶14 — Title screening overview and total records]}
% ↑ Delete this line and write your paragraph here.

% -----------------------------------------------------------------------
\subsubsection{Layer~1 — Hard Exclusions}
% -----------------------------------------------------------------------

%% WRITE ¶15 — Hard exclusions (~2 sentences):
%   • Records published before 2021 were excluded automatically.
%   • No minimum citation threshold was applied (relevance, not popularity,
%     drives inclusion at this stage — consistent with Wohlin 2014).

\bigskip\noindent\textit{[¶15 — Layer 1: year-based hard exclusions]}
% ↑ Delete this line and write your paragraph here.

% -----------------------------------------------------------------------
\subsubsection{Layer~2 — Keyword Scoring}
% -----------------------------------------------------------------------

%% WRITE ¶16 — Keyword term-set scoring (~3–4 sentences):
%   • Titles scored against 4 term sets: PEFT/adapters, LLM/transformers,
%     systems/serving, distributed/P2P/federated.
%   • Decision rules:
%       ≥ 2 term sets matched → INCLUDE
%       Exactly 1 non-LLM term set → REVIEW queue
%       LLM-only match → EXCLUDE
%       0 matches → EXCLUDE
%   • Note: exact case-insensitive substring matching, term-set version 1.0.
%   • State reproducibility: term set documented in screening audit log.

\bigskip\noindent\textit{[¶16 — Layer 2: keyword term-set scoring and decision rules]}
% ↑ Delete this line and write your paragraph here.

% -----------------------------------------------------------------------
\subsubsection{Layer~3 — LLM Triage of the REVIEW Queue}
% -----------------------------------------------------------------------

%% WRITE ¶17 — LLM triage (~3–4 sentences):
%   REQUIRED BY: Xiao & Watson (2019) — document tool/model identifiers;
%                all decisions must be auditable.
%   • 21 REVIEW-queue records submitted to Claude Haiku
%     (\texttt{claude-haiku-4-5-20251001}; Anthropic) for title- and
%     venue-level classification.
%   • Model returned: INCLUDE, EXCLUDE, or UNCERTAIN.
%   • UNCERTAIN records retained for manual inspection.
%   • State: AI decisions logged per record alongside keyword scores in
%     \texttt{log\_screening\_*.json} (SHA-256 of input file retained).

\bigskip\noindent\textit{[¶17 — Layer 3: LLM triage of REVIEW queue]}
% ↑ Delete this line and write your paragraph here.

\paragraph{Screening outcomes.}

% Fill in from log_screening_*.json after final pipeline run
\begin{table}[htbp]
  \centering
  \caption{Title screening outcomes across all 1,008 records.}
  \label{tab:screening}
  \begin{tabular}{lr}
    \toprule
    \textbf{Decision} & \textbf{\textit{N}} \\
    \midrule
    \textsc{Include} (automatic, $\geq 2$ term sets) & 162 \\
    \textsc{Review} (LLM triage queue)               &  21 \\
    \quad $\hookrightarrow$ \textsc{Include} after LLM triage & \TODO{$n$} \\
    \quad $\hookrightarrow$ \textsc{Exclude} after LLM triage & \TODO{$n$} \\
    \quad $\hookrightarrow$ \textsc{Uncertain} (manual)       & \TODO{$n$} \\
    \textsc{Exclude} (hard)                          & 825 \\
    \midrule
    \textbf{Total screened}                          & \textbf{1,008} \\
    \bottomrule
  \end{tabular}
\end{table}

%% WRITE ¶18 — Inter-rater reliability (~3–4 sentences):
%   REQUIRED BY: Xiao & Watson (2019) Step 4 — TWO reviewers must screen
%                independently; discrepancies resolved by discussion or
%                a third party; Cohen's kappa should be reported.
%   THIS IS A REQUIRED PARAGRAPH — do not omit it.
%   Options (pick one and state it clearly):
%   Option A — if a second reviewer participated:
%       Describe the second reviewer, the sample screened independently,
%       and the Cohen's kappa / percentage agreement achieved.
%   Option B — single reviewer (most common in thesis work):
%       Acknowledge explicitly that screening was performed by a single
%       reviewer (a recognised limitation of the SLR).
%       Describe the mitigations in place:
%         (i)  The LLM triage in Layer 3 provides a documented,
%              independent second perspective on the REVIEW queue.
%         (ii) The pre-validated G1–G6 corpora, assembled through a
%              separate search process, serve as a cross-validation
%              mechanism for thematic coverage.
%         (iii) Full decision audit logs enable post-hoc inter-rater
%              checks on any subset of records.
%   Cite: \citep{XiaoWatson2019, KitchenhamCharters2007}

\bigskip\noindent\textit{[¶18 — Inter-rater reliability (REQUIRED — acknowledge or report)]}
% ↑ Delete this line and write your paragraph here.


% =============================================================================
\subsection{Metadata Enrichment and Quality Assessment}
\label{sec:methodology:enrichment}
% =============================================================================

% REQUIRED BY: Xiao & Watson (2019) Step 5 — quality assessment with a
%              pre-specified instrument; results reported.
%              Xiao & Watson (2019) Step 6 — data extraction form applied
%              to each included study.
%              Wohlin (2014) §3.4 — data extraction after inclusion.

%% WRITE ¶19 — Enrichment pipeline (~4–5 sentences):
%   • All records passing title screening enriched with abstracts and
%     keyword metadata via two sources:
%       1. Semantic Scholar batch API — abstracts and field-of-study
%          categories for all records with a SS paper ID.
%       2. OpenAlex — keyword topics for records still lacking metadata
%          after the SS pass, via DOI lookup.
%   • Per-row traceability: the \texttt{enrichment\_sources} column
%     records which sources contributed data for each paper
%     (values: \texttt{acl}, \texttt{semantic\_scholar}, \texttt{openalex}).
%   • After enrichment, a domain-relevance filter removed off-topic records
%     and deprioritised low-citation arXiv preprints lacking a Tier~1
%     signal (moved to a separate audit file for potential rescue).
%   • Cite: \citep{SemanticScholar, XiaoWatson2019}

\bigskip\noindent\textit{[¶19 — Enrichment pipeline: SS + OpenAlex + traceability column]}
% ↑ Delete this line and write your paragraph here.

%% WRITE ¶20 — Quality assessment instrument (~3–4 sentences):
%   REQUIRED BY: Xiao & Watson (2019) Step 5 — pre-specified quality
%                assessment instrument; results reported; quality informs
%                synthesis weighting.
%   THIS IS A REQUIRED PARAGRAPH.
%   • State that venue quality was assigned to each retained record using
%     a four-level taxonomy (Table~\ref{tab:venue-quality}).
%   • Describe the taxonomy:
%       top_venue    — CORE A* / Scopus Q1 (ICML, ICLR, NeurIPS, ACL,
%                      EMNLP, MLSys, Scopus Q1 journals)
%       peer_reviewed — other peer-reviewed conference or journal
%       preprint     — arXiv with no published venue
%       unknown      — venue missing or unresolved
%   • State how quality is used: top-tier peer-reviewed works are
%     prioritised in synthesis; preprints are included where no published
%     version exists (consistent with Tier 1 priority criterion).
%   • Note the instrument was pre-specified in the protocol.
%   • Cite: \citep{XiaoWatson2019}

\bigskip\noindent\textit{[¶20 — Quality assessment: venue quality taxonomy]}
% ↑ Delete this line and write your paragraph here.

\paragraph{Tier classification.}

%% WRITE ¶21 — Tier classification (~3–4 sentences):
%   • Each retained record classified into one of three relevance tiers
%     based on combined title, abstract, and keyword signal
%     (Table~\ref{tab:tiers}).
%   • Connect to Keshav (2007) three-pass model:
%       Tier 1 → full reads
%       Tier 2 → selective reads
%       Tier 3 → abstract-level assessment
%   • Cite: \citep{Keshav2007, XiaoWatson2019}

\bigskip\noindent\textit{[¶21 — Tier classification and reading-priority rationale]}
% ↑ Delete this line and write your paragraph here.

\begin{table}[htbp]
  \centering
  \caption{Relevance tier classification criteria and corpus distribution
           after enrichment and filtering.}
  \label{tab:tiers}
  \begin{tabular}{clrr}
    \toprule
    \textbf{Tier} & \textbf{Criteria} & \textbf{\textit{n}} & \textbf{\%} \\
    \midrule
    1 & PEFT + distributed/P2P + LLM serving (all three) & 108 & 20.9 \\
    2 & Two of the three core signals                    & 163 & 31.5 \\
    3 & Foundational or tangential (one signal)          & 246 & 47.6 \\
    \midrule
    \textbf{Total} & & \textbf{517} & \\
    \bottomrule
  \end{tabular}
\end{table}

\begin{table}[htbp]
  \centering
  \caption{Venue quality distribution of the 517-paper retained corpus.}
  \label{tab:venue-quality}
  \begin{tabular}{lrr}
    \toprule
    \textbf{Venue quality} & \textbf{\textit{n}} & \textbf{\%} \\
    \midrule
    Top-tier venue (CORE A* / Scopus Q1)       &  66 & 12.8 \\
    Peer-reviewed (other conference / journal)  & 209 & 40.4 \\
    Preprint (arXiv, no published venue)        & 219 & 42.4 \\
    Unknown / missing venue                     &  23 &  4.4 \\
    \bottomrule
  \end{tabular}
\end{table}

%% WRITE ¶22 — Data extraction form (~2–3 sentences):
%   REQUIRED BY: Xiao & Watson (2019) Step 6 — a structured extraction
%                form must be described; fields listed or made available.
%   • State the fields extracted per record (the columns of the enriched
%     CSV constitute the extraction schema):
%       title, authors, year, venue, DOI, SS paper ID, abstract,
%       keywords, enrichment_sources, venue_quality, tier, source_engine
%   • Note that the schema was pre-specified in the protocol and applied
%     uniformly via the automated enrichment pipeline.
%   • Cite: \citep{XiaoWatson2019}

\bigskip\noindent\textit{[¶22 — Data extraction form and fields]}
% ↑ Delete this line and write your paragraph here.


% =============================================================================
\subsection{Abstract-Level Review}
\label{sec:methodology:abstract-review}
% =============================================================================

% REQUIRED BY: Xiao & Watson (2019) Step 4 — abstract screening before
%              full text; Step 5 — quality and relevance confirmed.
%              Wohlin (2014) §3.2.3 — final inclusion based on full paper
%              (abstract review here is the eligibility → inclusion step).
%              Keshav (2007) — first-pass reading model.

%% WRITE ¶23 — Abstract review purpose and process (~4–5 sentences):
%   • State purpose: move from eligibility to inclusion (PRISMA stage 4).
%   • Ground in Keshav (2007): first pass (title, abstract, headings,
%     five Cs) decides whether a paper warrants further reading.
%   • Describe the AI-advisory mechanism:
%       Each abstract accompanied by a Claude Sonnet suggestion
%       (\texttt{claude-sonnet-4-6}; Anthropic): one-sentence summary,
%       one-sentence relevance note, KEEP / SKIP / DEFER recommendation.
%   • CRITICAL: state clearly that the AI suggestion was STRICTLY ADVISORY;
%     ALL final decisions were made by the human reviewer.
%     (Required by Xiao & Watson: human oversight for validity.)
%   • State the export step: KEEP + DEFER → Zotero via RIS format.
%   • Cite: \citep{Keshav2007, XiaoWatson2019}

\bigskip\noindent\textit{[¶23 — Abstract review: process, AI advisory role, human decision]}
% ↑ Delete this line and write your paragraph here.

% TODO: fill in final counts after abstract review session
\begin{table}[htbp]
  \centering
  \caption{Abstract-level review outcomes (to be completed after
           \texttt{python -m app.abstract\_review --llm}).}
  \label{tab:abstract-review}
  \begin{tabular}{lr}
    \toprule
    \textbf{Decision} & \textbf{\textit{N}} \\
    \midrule
    KEEP  (Zotero corpus for full-text reading) & \TODO{$n$} \\
    SKIP  (excluded at abstract level)          & \TODO{$n$} \\
    DEFER (closer inspection required)          & \TODO{$n$} \\
    \midrule
    \textbf{Total reviewed}                     & \textbf{517} \\
    \bottomrule
  \end{tabular}
\end{table}

%% WRITE ¶24 — Reproducibility of abstract review (~2 sentences):
%   • Both the AI suggestion and the human decision are recorded per record
%     in \texttt{08\_abstract\_reviewed\_<date>.csv}.
%   • This supports inter-rater reproducibility checks on any subset of
%     SKIP decisions (the AI column serves as a first-pass surrogate for
%     a second reviewer — Xiao & Watson recommendation).
%   • Cite: \citep{XiaoWatson2019, KitchenhamCharters2007}

\bigskip\noindent\textit{[¶24 — Reproducibility of abstract review decisions]}
% ↑ Delete this line and write your paragraph here.


% =============================================================================
\subsection{PRISMA Flow and Reproducibility}
\label{sec:methodology:prisma}
% =============================================================================

% REQUIRED BY: Page et al. (2021)   — PRISMA 2020 flow diagram (mandatory)
%              Xiao & Watson (2019) Step 8 — report with flow diagram;
%              document search date, string, all decisions.
%              Xiao & Watson: "process must be reported in sufficient detail
%              to allow other researchers to follow the same steps and
%              arrive at the same results."

%% WRITE ¶25 — PRISMA flow summary (~3–4 sentences):
%   • Reference Figure~\ref{fig:prisma} for the complete PRISMA 2020 flow.
%   • Briefly narrate the four stages:
%       Identification  — 1,150 raw records retrieved across backward
%                         and forward snowballing.
%       Screening       — after deduplication: 41 unique new snowballed
%                         records + 444 pre-validated = 1,008 total;
%                         825 excluded at title screen.
%       Eligibility     — 517 records after enrichment and relevance filter.
%       Inclusion       — \TODO{$n$} KEEP + DEFER after abstract review.
%   • Cite: \citep{Page2021}

\bigskip\noindent\textit{[¶25 — PRISMA flow narrative (Identification → Inclusion)]}
% ↑ Delete this line and write your paragraph here.

%% WRITE ¶26 — Reproducibility statement (~2–3 sentences):
%   REQUIRED BY: Xiao & Watson (2019) — explicit reproducibility standard.
%   • State that the pipeline, term-set versions, LLM model identifiers,
%     SHA-256 checksums of all intermediate files, and per-seed retrieval
%     logs are archived in the project repository.
%   • State that full re-execution from any stage is possible.
%   • Cite: \citep{Page2021, XiaoWatson2019}

\bigskip\noindent\textit{[¶26 — Reproducibility and archival statement]}
% ↑ Delete this line and write your paragraph here.

% ---- PRISMA flow figure ----
% Replace \fbox placeholder with \includegraphics once generated:
%   python -m app.prisma --date <DATE>
%   python main.py --visualise
% Produces: slr_engine/assets/figures/fig_slr1_prisma_flow.png
\begin{figure}[htbp]
  \centering
  % \includegraphics[width=\textwidth]{slr_engine/assets/figures/fig_slr1_prisma_flow.png}
  \fbox{\parbox{0.88\textwidth}{%
    \centering\small
    \textbf{PRISMA 2020 Flow Diagram}\\[8pt]
    %
    \textbf{Identification}\\
    Records retrieved via backward snowballing: \TODO{$n$}\\
    Records retrieved via forward snowballing: \TODO{$n$}\\
    \textit{Total raw records: 1,150}\\[6pt]
    $\downarrow$ Deduplication (DOI + SS paper ID)\\[6pt]
    %
    \textbf{Screening}\\
    Unique new snowballed records: 41
      \quad (duplicates removed: 1,109)\\
    Pre-validated records (G0--G6): 444\\
    \textit{Total records screened: 1,008}\\
    Records excluded at title screen: 825\\[6pt]
    $\downarrow$\\[6pt]
    %
    \textbf{Eligibility}\\
    Records after enrichment \& relevance filter: 517\\
    \quad of which Tier~1: 108 \quad Tier~2: 163 \quad Tier~3: 246\\
    Records deprioritised (low-citation preprint): \TODO{$n$}\\
    Records excluded (off-topic / malformed): \TODO{$n$}\\[6pt]
    $\downarrow$\\[6pt]
    %
    \textbf{Inclusion}\\
    KEEP (full-text reading): \TODO{$n$}\\
    DEFER (closer inspection): \TODO{$n$}\\
    SKIP (excluded at abstract): \TODO{$n$}\\
    \textbf{Final corpus for full-text synthesis: \TODO{$n$}}
  }}
  \caption{PRISMA~2020 flow diagram.  Counts marked \TODO{\$n\$} are
           filled upon completion of the abstract review stage by
           re-running \texttt{python -m app.prisma} and
           \texttt{python main.py --visualise}.}
  \label{fig:prisma}
\end{figure}


% =============================================================================
%  Bibliography entries referenced in this section
%  (add to your main .bib file if not already present)
% =============================================================================

% @article{XiaoWatson2019,
%   author  = {Xiao, Yu and Watson, Maria},
%   title   = {Guidance on Conducting a Systematic Literature Review},
%   journal = {Journal of Planning Education and Research},
%   volume  = {39},
%   number  = {1},
%   pages   = {93--112},
%   year    = {2019},
%   doi     = {10.1177/0739456X17723971}
% }
%
% @inproceedings{Wohlin2014,
%   author    = {Wohlin, Claes},
%   title     = {Guidelines for Snowballing in Systematic Literature Studies
%                and a Replication in Software Engineering},
%   booktitle = {Proceedings of the 18th International Conference on Evaluation
%                and Assessment in Software Engineering (EASE~'14)},
%   year      = {2014},
%   doi       = {10.1145/2601248.2601268}
% }
%
% @article{Page2021,
%   author  = {Page, Matthew J. and others},
%   title   = {{PRISMA} 2020 explanation and elaboration: updated guidance
%              and exemplars for reporting systematic reviews},
%   journal = {BMJ},
%   volume  = {372},
%   year    = {2021},
%   doi     = {10.1136/bmj.n160}
% }
%
% @techreport{KitchenhamCharters2007,
%   author      = {Kitchenham, Barbara and Charters, Stuart},
%   title       = {Guidelines for performing systematic literature reviews
%                  in software engineering},
%   institution = {Keele University},
%   year        = {2007},
%   number      = {EBSE-2007-01}
% }
%
% @article{Keshav2007,
%   author  = {Keshav, Srinivasan},
%   title   = {How to read a paper},
%   journal = {ACM SIGCOMM Computer Communication Review},
%   volume  = {37},
%   number  = {3},
%   pages   = {83--84},
%   year    = {2007}
% }
%
% @misc{SemanticScholar,
%   author       = {{Semantic Scholar}},
%   title        = {Semantic Scholar Academic Graph {API}},
%   howpublished = {\url{https://api.semanticscholar.org}},
%   year         = {2024}
% }
%
% -- Seed paper entries assumed already in .bib:
%    Houlsby2019, Hu2022, Pfeiffer2020hub, Pfeiffer2021fusion,
%    Sajina2024, Borzunov2023, Han2024, Sheng2024
